% Figure: the Langmuir adsorption isotherm for the strongly-adsorbed
% component on the PSA sorbent, with the working capacity marked as the swing
% in loading between the high adsorption pressure and the low desorption
% pressure. The favourable (concave-down) shape is why the swing is taken
% across a partial-pressure interval, not from zero.
\begin{figure}[htbp]
\centering
\begin{tikzpicture}
\begin{axis}[
  width=11cm, height=6.6cm,
  xlabel={partial pressure of adsorbate $p$ (bar)},
  ylabel={loading $q$ (mol/kg)},
  xmin=0, xmax=6, ymin=0, ymax=3.2,
  legend pos=south east,
  legend cell align=left,
  grid=both, grid style={enggray!18},
  samples=120, domain=0:6,
]
% Langmuir isotherm q = q_m b p /(1 + b p), q_m = 3.0 mol/kg, b = 0.9 /bar
\addplot[engblue, very thick] {3.0*0.9*x/(1+0.9*x)};
\addlegendentry{Langmuir isotherm}
% desorption point at p_L = 1 bar : q = 3*0.9/(1.9) = 1.421
\draw[enggray, dashed] (axis cs:1,0) -- (axis cs:1,1.421);
\draw[enggray, dashed] (axis cs:0,1.421) -- (axis cs:1,1.421);
\addplot[only marks, mark=*, mark size=1.6pt, black] coordinates {(1,1.421)};
\node[font=\scriptsize, anchor=north west] at (axis cs:1.05,1.30) {$q_L$ at $p_L$};
% adsorption point at p_H = 5 bar : q = 3*4.5/(5.5) = 2.455
\draw[enggray, dashed] (axis cs:5,0) -- (axis cs:5,2.455);
\draw[enggray, dashed] (axis cs:0,2.455) -- (axis cs:5,2.455);
\addplot[only marks, mark=*, mark size=1.6pt, black] coordinates {(5,2.455)};
\node[font=\scriptsize, anchor=south west] at (axis cs:4.2,2.50) {$q_H$ at $p_H$};
% working-capacity swing arrow on the loading axis
\draw[engred, very thick, {Latex}-{Latex}] (axis cs:0.35,1.421) -- (axis cs:0.35,2.455);
\node[font=\scriptsize, engred, anchor=west, align=left] at (axis cs:0.45,1.94)
  {working\\ capacity\\ $\Delta q$};
\end{axis}
\end{tikzpicture}
\caption[Langmuir isotherm and the PSA working capacity]{The adsorption
isotherm of the strongly-held component fixes the pressure-swing working
capacity. At the high adsorption pressure $p_H$ the sorbent loads to
$q_H$; at the low desorption pressure $p_L$ it gives back down to $q_L$,
and the swing $\Delta q = q_H - q_L$ is the amount of contaminant each
kilogram of sorbent removes per cycle. The favourable concave-down shape
of the Langmuir isotherm is a mixed blessing for a pressure swing: it
gives a high loading at modest pressure, but it also flattens at high
pressure, so most of the swing comes from the lower part of the curve and
pushing $p_H$ higher buys little extra capacity. Reading the swing off
the isotherm at the two operating pressures is the first sizing move of
any PSA design.}
\label{fig:vol04ch07:psa-isotherm}
\end{figure}
